% GENERATED FILE. Edit canonical lessons, then run npm run generate:books.
% canonical-source-hash: fnv1a64:0e63e55ddb11cf51
% canonical-lessons: FA-C335-arusi, FA-C335-arus, FA-C335-hanabandan, FA-C335-mehriye, FA-C335-jahiziye

\chapter{The wedding}
\label{ch:fa-the-wedding}
\hlchaptermodality{335}

\begin{chapteropening}
\textbf{By the end of this chapter:} \emph{I can talk about a wedding, the bride, the henna night, the bride's marriage gift and her trousseau.}

Complete the last lesson of chapter 335: I can talk about a wedding, the bride, the henna night, the bride's marriage gift and her trousseau.
\end{chapteropening}

\section[\texorpdfstring{\fa{عروسی} ('arusi)}{عروسی ('arusi)}]{\fa{عروسی} ('arusi) — a wedding}

\label{lesson:FA-C335-arusi}



\begin{quote}
Before the new one: say the Persian for an engagement, then the Persian for the marriage contract.
\end{quote}

\subsection*{You'll want to know: \fa{عروسی}}
\textbf{\fa{عروسی}} (\emph{'arusi}) — \textquotedblleft{}a wedding\textquotedblright{}.

\begin{quote}
\fa{عروسی} \fa{دوستم} \fa{تا} \fa{نیمه} \fa{شب} \fa{طول} \fa{کشید}.
\end{quote}

\begin{quote}
\emph{'Arusi-ye dustam tâ nime shab tul keshid.}
\end{quote}

\begin{quote}
My friend's wedding went on until midnight.
\end{quote}

Built on \emph{'arus} \textquotedblleft{}bride\textquotedblright{}. In speech \emph{'arusi kardan} is \textquotedblleft{}to get married\textquotedblright{}.

\begin{cousinweb}
\emph{'arusi} adds the Persian suffix \emph{-i} to \emph{'arus} \textquotedblleft{}bride\textquotedblright{}, which is borrowed from Arabic \emph{'arūs} (EtymDB).
\end{cousinweb}

\begin{grammarlens}[title={What you've built}]
One more word for telling, explaining and describing at B1.
\end{grammarlens}

\subsection*{Your turn}
\begin{itemize}
\raggedright
  \item \emph{Say it:} \emph{'arusi}
  \item \emph{Say it:} \emph{'arusi}, once more
  \item \emph{Say it:} a sentence of your own with \emph{'arusi}
  \item \emph{Recall:} say the Persian for an engagement, then the Persian for the marriage contract, then say \emph{'arusi} again
\end{itemize}

\begin{tcolorbox}[breakable,colback=teal!4,colframe=teal!35!black,title={Before you move on}]
What does \fa{عروسی} mean? (\textquotedblleft{}A wedding\textquotedblright{}.) Say it once more.
\end{tcolorbox}
\section[\texorpdfstring{\fa{عروس} ('arus)}{عروس ('arus)}]{\fa{عروس} ('arus) — a bride; a daughter-in-law}

\label{lesson:FA-C335-arus}



\begin{quote}
Before the new one: say the Persian for the marriage contract, then the Persian for a wedding.
\end{quote}

\subsection*{You'll want to know: \fa{عروس}}
\textbf{\fa{عروس}} (\emph{'arus}) — \textquotedblleft{}a bride; a daughter-in-law\textquotedblright{}.

\begin{quote}
\fa{عروس} \fa{لباس} \fa{سفید} \fa{پوشیده} \fa{بود}.
\end{quote}

\begin{quote}
\emph{'Arus lebâs-e sefid pushide bud.}
\end{quote}

\begin{quote}
The bride was wearing a white dress.
\end{quote}

It also means \textquotedblleft{}daughter-in-law\textquotedblright{}. The groom is \emph{dâmâd}.

\begin{cousinweb}
\emph{'arus} is borrowed from Arabic \emph{'arūs} (EtymDB), which can name either the bride or the groom; in Persian it is the bride.
\end{cousinweb}

\begin{grammarlens}[title={What you've built}]
One more word for telling, explaining and describing at B1.
\end{grammarlens}

\subsection*{Your turn}
\begin{itemize}
\raggedright
  \item \emph{Say it:} \emph{'arus}
  \item \emph{Say it:} \emph{'arus}, once more
  \item \emph{Say it:} a sentence of your own with \emph{'arus}
  \item \emph{Recall:} say the Persian for the marriage contract, then the Persian for a wedding, then say \emph{'arus} again
\end{itemize}

\begin{tcolorbox}[breakable,colback=teal!4,colframe=teal!35!black,title={Before you move on}]
What does \fa{عروس} mean? (\textquotedblleft{}A bride; a daughter-in-law\textquotedblright{}.) Say it once more.
\end{tcolorbox}
\section[\texorpdfstring{\fa{حنابندان} (hanâbandân)}{حنابندان (hanâbandân)}]{\fa{حنابندان} (hanâbandân) — a henna night}

\label{lesson:FA-C335-hanabandan}



\begin{quote}
Before the new one: say the Persian for a wedding, then the Persian for a bride.
\end{quote}

\subsection*{You'll want to know: \fa{حنابندان}}
\textbf{\fa{حنابندان}} (\emph{hanâbandân}) — \textquotedblleft{}a henna night\textquotedblright{}.

\begin{quote}
\fa{شب} \fa{حنابندان} \fa{زنها} \fa{آواز} \fa{میخواندند}.
\end{quote}

\begin{quote}
\emph{Shab-e hanâbandân zan-hâ âvâz mikhândand.}
\end{quote}

\begin{quote}
On the henna night the women sang.
\end{quote}

The night before the wedding, when henna is put on the bride's hands. \emph{hanâ} \textquotedblleft{}henna\textquotedblright{} + \emph{band} \textquotedblleft{}binding\textquotedblright{}.

\begin{cousinweb}
\emph{hanâbandân} joins \emph{hanâ} \textquotedblleft{}henna\textquotedblright{}, borrowed from Arabic \emph{ḥinnā'} (EtymDB), and the native Persian \emph{band-} \textquotedblleft{}to bind\textquotedblright{} with the suffix \emph{-ân}. EtymDB lists no entry.
\end{cousinweb}

\begin{grammarlens}[title={What you've built}]
One more word for telling, explaining and describing at B1.
\end{grammarlens}

\subsection*{Your turn}
\begin{itemize}
\raggedright
  \item \emph{Say it:} \emph{hanâbandân}
  \item \emph{Say it:} \emph{hanâbandân}, once more
  \item \emph{Say it:} a sentence of your own with \emph{hanâbandân}
  \item \emph{Recall:} say the Persian for a wedding, then the Persian for a bride, then say \emph{hanâbandân} again
\end{itemize}

\begin{tcolorbox}[breakable,colback=teal!4,colframe=teal!35!black,title={Before you move on}]
What does \fa{حنابندان} mean? (\textquotedblleft{}A henna night\textquotedblright{}.) Say it once more.
\end{tcolorbox}
\section[\texorpdfstring{\fa{مهریه} (mehriye)}{مهریه (mehriye)}]{\fa{مهریه} (mehriye) — a mahr, the bride's marriage gift}

\label{lesson:FA-C335-mehriye}



\begin{quote}
Before the new one: say the Persian for a bride, then the Persian for a henna night.
\end{quote}

\subsection*{You'll want to know: \fa{مهریه}}
\textbf{\fa{مهریه}} (\emph{mehriye}) — \textquotedblleft{}a mahr, the bride's marriage gift\textquotedblright{}.

\begin{quote}
\fa{مهریه} \fa{او} \fa{صد} \fa{سکه} \fa{طلا} \fa{بود}.
\end{quote}

\begin{quote}
\emph{Mehriye-ye u sad sekke-ye talâ bud.}
\end{quote}

\begin{quote}
Her mahr was a hundred gold coins.
\end{quote}

Promised by the groom in the marriage contract and owed to the wife.

\begin{cousinweb}
\emph{mehriye} is built on \emph{mehr}, the marriage gift, with the ending \emph{-iye}. \emph{mehr} is usually traced to Arabic \emph{mahr}, \textquotedblleft{}marriage portion\textquotedblright{} (EtymDB lists the Arabic word), though EtymDB also links the Persian word to an Old Persian form, so its origin is not fully settled.
\end{cousinweb}

\begin{grammarlens}[title={What you've built}]
One more word for telling, explaining and describing at B1.
\end{grammarlens}

\subsection*{Your turn}
\begin{itemize}
\raggedright
  \item \emph{Say it:} \emph{mehriye}
  \item \emph{Say it:} \emph{mehriye}, once more
  \item \emph{Say it:} a sentence of your own with \emph{mehriye}
  \item \emph{Recall:} say the Persian for a bride, then the Persian for a henna night, then say \emph{mehriye} again
\end{itemize}

\begin{tcolorbox}[breakable,colback=teal!4,colframe=teal!35!black,title={Before you move on}]
What does \fa{مهریه} mean? (\textquotedblleft{}A mahr, the bride's marriage gift\textquotedblright{}.) Say it once more.
\end{tcolorbox}
\section[\texorpdfstring{\fa{جهیزیه} (jahiziye)}{جهیزیه (jahiziye)}]{\fa{جهیزیه} (jahiziye) — a trousseau, the bride's household goods}

\label{lesson:FA-C335-jahiziye}



\begin{quote}
Before the new one: say the Persian for a henna night, then the Persian for a mahr.
\end{quote}

\subsection*{You'll want to know: \fa{جهیزیه}}
\textbf{\fa{جهیزیه}} (\emph{jahiziye}) — \textquotedblleft{}a trousseau, the bride's household goods\textquotedblright{}.

\begin{quote}
\fa{مادرش} \fa{جهیزیه} \fa{را} \fa{از} \fa{یک} \fa{سال} \fa{پیش} \fa{خریده} \fa{بود}.
\end{quote}

\begin{quote}
\emph{Mâdarash jahiziye râ az yek sâl pish kharide bud.}
\end{quote}

\begin{quote}
Her mother had bought the trousseau a year before.
\end{quote}

The furniture and household things that the bride's family provides.

\begin{cousinweb}
\emph{jahiziye} is built on \emph{jahiz} \textquotedblleft{}a trousseau\textquotedblright{}, which EtymDB derives from Arabic \emph{jihāz}, \textquotedblleft{}equipment, outfit\textquotedblright{}, with the ending \emph{-iye}.
\end{cousinweb}

\begin{grammarlens}[title={What you've built}]
One more word for telling, explaining and describing at B1.
\end{grammarlens}

\subsection*{Your turn}
\begin{itemize}
\raggedright
  \item \emph{Say it:} \emph{jahiziye}
  \item \emph{Say it:} \emph{jahiziye}, once more
  \item \emph{Say it:} a sentence of your own with \emph{jahiziye}
  \item \emph{Recall:} say the Persian for a henna night, then the Persian for a mahr, then say \emph{jahiziye} again
\end{itemize}

\begin{tcolorbox}[breakable,colback=teal!4,colframe=teal!35!black,title={Before you move on}]
What does \fa{جهیزیه} mean? (\textquotedblleft{}A trousseau, the bride's household goods\textquotedblright{}.) Say it once more.
\end{tcolorbox}
